\chapter{Introduction}
\label{chap:intro}

\section{Background and Motivation}

One of the workhorses of modern hospital imaging is the abdominal CT scan. Tumour diagnosis, staging, surgical planning, radiotherapy target definition --- all of these processes rely at some stage on understanding precisely where each organ is located in the volume. Without that spatial context, downstream decisions --- how much dose to give, which margin to cut out, which structure to avoid --- are just educated guesses. This is a large problem: a busy radiology department can receive dozens of abdominal CTs a day, and each one ideally should result in a pixel-wise map of the organs before treatment planning can even get started.

Manual organ delineation is an extremely tedious and time-consuming process. For a single CT volume consisting of $200$--$400$ axial slices, requiring delineation of 13 structures, a trained radiologist would usually spend between one and three hours producing an accurate segmentation. When scaled across the entire scan volume of a large hospital this bottleneck becomes apparent --- leading to delays in treatment planning, inter-observer variation, and higher-value diagnostic tasks being neglected. Automated segmentation reduces this burden, turning hours of manual work into seconds of computation.

Deep convolutional architectures have radically changed the game for segmentation. The encoder--decoder framework popularised by U-Net~\cite{ronneberger2015} has been extended in many directions: by including attention mechanisms, by handling volumetric data, and by developing self-supervised pretraining algorithms. However, there remain some notable gaps. 3D approaches are computationally intensive, transformer approaches are data-hungry, and many published results still optimistically report per-slice evaluation metrics instead of per-volume metrics. This project aims to address these shortcomings specifically, through a careful implementation of a 2D Attention U-Net with a rigorous per-volume evaluation protocol.

\begin{figure}[h!]
\centering
\begin{minipage}{0.97\textwidth}
  \centering
  \includegraphics[width=\textwidth]{figures/fig_pipeline}
  \caption{End-to-end pipeline: raw CT volume $\rightarrow$ HU windowing $[-175, 250]$
  $\rightarrow$ per-slice resize to $256\times256$ $\rightarrow$ Attention U-Net inference
  $\rightarrow$ multi-organ segmentation mask (13 organ classes).}
  \label{fig:pipeline}
\end{minipage}
\end{figure}

\FloatBarrier

\section{Problem Statement}

Our problem is \emph{multi-organ semantic segmentation}: for every axial CT slice, assign every pixel a label from 14 options (background plus 13 abdominal organs). Three specific difficulties make this daunting. First, class ratios are extremely skewed --- the combined foreground voxels for the pancreas and adrenal glands are under 1\%. Second, shape distributions are highly variable across patients, especially for the stomach and duodenum. Third, 2D slice-by-slice processing is necessary due to GPU memory constraints on freely available hardware.

A second caveat that only became obvious during our initial experiments is evaluation integrity. We observed that averaging Dice scores computed slice-by-slice gives misleadingly high average scores: the underlying problem is empty slices. When the prediction and the ground-truth are both all-zero, a na\"\i{}ve implementation with a smoothing term computes the Dice as 1.0. Though a handful of such slices inflates the average score for each volume, it actually obscures the true segmentation performance. We addressed this on a per-volume basis, using the protocol described in this report.

\section{Objectives}

\begin{itemize}
    \item Implement a 2D Attention U-Net with attention gates on all skip connections for
          multi-organ abdominal CT segmentation.
    \item Train and evaluate on the FLARE22 benchmark using a per-volume evaluation protocol
          consistent with the challenge's own scoring conventions.
    \item Quantify the contribution of attention gates, Test-Time Augmentation, and cosine
          learning-rate scheduling through a controlled ablation study.
    \item Show competitive Dice and HD95 results against published nnU-Net 3D baselines,
          while using only 2D processing and a single T4 GPU.
\end{itemize}

\section{Contributions}

The principal contributions of this project are summarised in Table~\ref{tab:contributions}.

\begin{table}[h!]
\centering
\caption{Summary of contributions.}
\label{tab:contributions}
\begin{tabular}{p{3.5cm} p{7cm} p{2.5cm}}
\toprule
\textbf{Contribution} & \textbf{Description} & \textbf{Chapter} \\
\midrule
Attention U-Net implementation &
  Full PyTorch implementation with attention gates on all four skip connections, Dropout2d
  bottleneck, and AMP-enabled training. &
  Chapter~\ref{chap:methodology} \\[6pt]
Per-volume evaluation fix &
  Identified and corrected the per-slice Dice inflation bug; established a per-volume evaluation
  loop consistent with FLARE22 conventions. &
  Chapter~\ref{chap:eval} \\[6pt]
TTA ensemble &
  Four-way horizontal/vertical flip ensemble with softmax averaging, adding $+$0.06 Dice at zero
  extra training cost. &
  Chapter~\ref{chap:eval} \\[6pt]
Ablation study &
  Controlled comparison across Runs~014--017 isolating the effect of data volume, attention gates,
  cosine LR, and TTA. &
  Chapter~\ref{chap:ablation} \\[6pt]
Competitive results on tight hardware &
  Mean Dice 0.8769, HD95 7.59~mm on 13 organs; numerically exceeding the nnU-Net 3D baseline
  on 7 of 13 organs using only 2D slices and a single T4 GPU. &
  Chapter~\ref{chap:results} \\
\bottomrule
\end{tabular}
\end{table}

\section{Report Organisation}

Chapter~\ref{chap:lit} and Chapter~\ref{chap:data} provide an overview of the related work on segmentation of medical images, attention mechanisms, and the FLARE22 benchmark respectively. Chapter~\ref{chap:methodology} provides a detailed explanation of the Attention U-Net architecture and the loss function. Chapter~\ref{chap:train} introduces the training configuration, while Chapter~\ref{chap:eval} details our evaluation protocol and justifies our choice. Chapter~\ref{chap:results} presents the quantitative results, and Chapter~\ref{chap:ablation} covers the ablation study and interpretation of the results. The conclusion and future directions can be found in Chapters~\ref{chap:conc} and \ref{chap:future}.